\section{Counting, unary periodicity, and balanced parentheses}
\label{sec:counting}

This section proves numerical consequences of the finite constructions.
Formal power-series notation below always abbreviates a statement about
every supplied finite truncation. No completed infinite sequence is used
as an object of the arithmetic.

\begin{theorem}[Counting accepted words by a finite matrix]
\label{count:matrix}
For a DFA on $r\ge1$ states over $q$ letters, let
\[
 M_{ij}=\#\{a<q:\delta(i,a)=j\},\qquad
 u_i=[i=q_0],\qquad f_i=[i\text{ is accepting}].
\]
The number $c_n$ of accepted words of length $n$ is
$c_n=uM^nf$. All entries, powers, and counts are primitive recursive,
uniformly in the finite DFA and $n$, and $0\le c_n\le q^n$.
\end{theorem}
\begin{proof}
Let $v_{n,j}$ count words of length $n$ ending at state $j$.
At length zero its value is $u_j$. Every word of length $n+1$ has a
unique length-$n$ prefix and a final letter. Partition those prefixes
by their reached state and then count the letters leading to $j$:
\[
 v_{n+1,j}=\sum_{i<r}v_{n,i}M_{ij}.
\]
The partition is a finite bijection, so its equality follows by induction
on the finite enumeration, with equality of running counts as invariant.
Inducting on $n$ gives $v_n=uM^n$. Summing over accepting states proves
the formula. At each state the row sum of $M$ is $q$; thus
$\sum_jv_{n,j}=q^n$ by the same recurrence, and the accepting sum is
bounded by the total. This argument counts words, not accepting paths
of a nondeterministic automaton, where multiplicities could be wrong.
\end{proof}

\begin{lemma}[Finite adjugate identity]\label{count:adjugate}
For a finite square matrix $A$ over integer polynomials,
\[
 A\operatorname{adj}(A)=\operatorname{adj}(A)A=\det(A)I.
\]
The polynomial coefficients on both sides are primitive-recursive finite
expressions, and their equality has an \ipra{} proof with signed integers
represented by pairs of natural numbers.
\end{lemma}
\begin{proof}
Use the convention that the determinant of the empty matrix is $1$.
Define the determinant by its finite permutation sum with the sign given
by the parity of the inversion count. For fixed row $i$, group the
permutations by the column used in that row. Removing that row and column
identifies the corresponding summand with the signed minor, and gives
Laplace expansion. The sign changes by $(-1)^{i+j}$: moving row $i$ and
column $j$ to the first positions uses $i+j$ adjacent exchanges, and
the remaining inversion parity is that of the smaller permutation.
Each grouping and deletion has a bounded inverse on the corresponding
finite permutation table.

For a diagonal entry of $A\operatorname{adj}(A)$, this is exactly the
expansion of $\det(A)$. For an off-diagonal entry, it is the determinant
of the matrix obtained by repeating one row in place of another. Pair
each permutation with the permutation exchanging the columns assigned
to those two row positions. The product terms agree, their signs differ,
and the involution has no fixed points, so the finite signed sum is zero.
Column expansion gives the other product. Arithmetic of integer pairs
reduces the signed identities to equalities after moving negative sums
to the other side. Polynomial multiplication is finite convolution,
so the same reasoning applies coefficient by coefficient. These are
bounded rearrangements of at most $r!$ terms; induction is only on the
number of terms and the polynomial coefficient position.
\end{proof}

\begin{theorem}[Rational counting series and an integer recurrence]
\label{count:rational}
From any $r$-state DFA one can primitive-recursively compute integer
polynomials
\[
 D(z)=\det(I-zM)=1+d_1z+\cdots+d_rz^r,\qquad
 P(z)=u\operatorname{adj}(I-zM)f,\quad\deg P\le r-1,
\]
such that, coefficientwise,
\[
 D(z)\sum_{n\ge0}c_nz^n=P(z).
\]
In particular, for every $n\ge r$,
\[
 c_n+d_1c_{n-1}+\cdots+d_rc_{n-r}=0.
\]
These assertions are \ipra{} theorems about finite coefficients.
\end{theorem}
\begin{proof}
Fix a truncation degree $N$ and work modulo $z^{N+1}$, a finite ring of
$(N+1)$ signed coefficient cells. The finite geometric sum
$S_N=I+zM+\cdots+z^NM^N$ satisfies
$(I-zM)S_N=I$ modulo $z^{N+1}$, by adjacent-term cancellation.
Multiplying by the adjugate and using Lemma~\ref{count:adjugate} gives
$D(z)S_N=\operatorname{adj}(I-zM)$ in that finite ring.
Multiply by $u$ and $f$ to obtain the displayed coefficient identity
through degree $N$. The constant determinant is $1$, since the constant
matrix is $I$; each minor has at most $r-1$ linear factors. The coefficient
of degree $n\ge r$ therefore gives the recurrence. Since $N$ was an
arbitrary numerical parameter, universal introduction proves every
coefficient equation. No convergence, eigenvalues, complex numbers,
or analytic division is needed. Conversely the equation with constant
denominator coefficient $1$ computes each next $c_n$ uniquely from
earlier coefficients, by primitive recursion with a finite history.
\end{proof}

\begin{theorem}[Unary regular languages are effectively ultimately periodic]
\label{count:unary}
For an $r$-state unary DFA there are computed numbers $i<r$ and
$1\le p\le r$ such that membership of $a^n$ equals membership of
$a^{n+p}$ for every $n\ge i$. Conversely, a unary primitive-recursive
language with supplied $i\ge0,p>0$ and this periodicity property is accepted
by an explicitly constructed DFA with $i+p$ states. The converse also
allows $i=0$.
\end{theorem}
\begin{proof}
Among the states reached at times $0,\ldots,r$, select equal states at
$i<j\le r$ and set $p=j-i$. Run composition gives equal states at
$i+t$ and $j+t$ for every $t$, by induction on $t$. Lookup of their
acceptance bits proves the first assertion. For the converse, make a
chain of states $0,\ldots,i+p-1$ whose final successor is state $i$.
The initial state is $0$, and state $j$ has the supplied language's
membership bit at $a^j$. For $n<i+p$ the reached state is $n$; afterwards
it is $i+((n-i)\bmod p)$. Repeated use of the given period, formalized
by induction on the quotient, proves the correct acceptance bit.
\end{proof}

\subsection{The one-type Dyck language}

Write an opening parenthesis as an up-step and a closing parenthesis as
a down-step. For a prefix $v$ let $U(v)$ and $D(v)$ be its counts of
these letters. Define the language $\mathcal D$ by
\[
 U(w)=D(w)\quad\text{and}\quad
 \forall j\le |w|\;D(w[0:j])\le U(w[0:j]).
\]
This is a primitive-recursive predicate. Height can be handled as a
signed pair, or by the two displayed nonnegative counts.

\begin{theorem}[Unique first return and an unambiguous grammar]
\label{count:dyck}
The language $\mathcal D$ is generated by
$S\to\eps\mid\mathtt{(}S\mathtt{)}S$.
Every nonempty Dyck word has a unique decomposition
$\mathtt{(}u\mathtt{)}v$ with $u,v\in\mathcal D$.
The matching closing position and the unique parse tree are
primitive-recursively computable.
\end{theorem}
\begin{proof}
A nonempty Dyck word starts with an opening parenthesis. Choose the first
positive prefix length at which $U=D$; it exists at the full length and
is found by bounded search. The last letter of this prefix is closing.
Before that point the height is at least one, so removing the first and
last steps gives a Dyck word $u$. The suffix $v$ is also Dyck because
the removed prefix has height zero. Conversely an opening step, any
Dyck word, a closing step, and another Dyck word never have negative
height and end at zero. Inside the first block the height is positive,
so this closing position is the unique first return.

Repeat the decomposition on the two strictly shorter words to build
the parse tree. Its nonempty internal nodes consume a disjoint pair
of parentheses, so there are at most $|w|/2$ such nodes. A work list
of subintervals and a loop of at most $|w|+1$ expansions implement the
recursion. For a word with $n$ pairs every parse has $2n+1$ variable
nodes and $2n$ terminal nodes: its variable nodes form a full binary
tree with $n$ nonempty expansions. This gives an explicit
primitive-recursive bound on all relevant tree codes. Uniqueness follows
by bounded induction on the word length, since the root split and then
both smaller parses are forced. The induction predicate checks equality
of all valid parses within the numerical bound. This proves
both grammar equivalence and unambiguity for this specific grammar.
\end{proof}

\begin{theorem}[Catalan recurrence and exact enumeration]
\label{count:catalan}
Let $C_n$ count Dyck words with $n$ opening and $n$ closing parentheses.
Then
\[
 C_0=1,\qquad C_{n+1}=\sum_{i=0}^{n}C_iC_{n-i},\qquad
 (n+1)C_n=\binom{2n}{n}.
\]
All counts and the bijections proving the identities are primitive
recursive, and the identities are provable in \ipra.
\end{theorem}
\begin{proof}
The empty word is the sole object for $n=0$. In the unique decomposition
of a nonempty word, $u$ has $i$ pairs and $v$ has $n-i$ pairs. The
decomposition and its inverse are a bijection, giving the convolution.

For the exact formula assume first $n\ge1$. There are $\binom{2n}{n}$
balanced words before imposing the prefix condition: choose the positions
of the opening steps. A bad balanced word has a first prefix of height
$-1$. Exchange opening and closing steps in this prefix. The resulting
word has $n+1$ opening and $n-1$ closing steps. Conversely any such word
has final height $2$ and therefore has a first prefix of height $1$;
exchange its steps. The result is balanced with a first prefix of height
$-1$, and the two operations are inverses. Existence of these first
crossings follows from the fact that each step changes height by one,
proved by a scan with the previous height as invariant. Thus the bad
words number $\binom{2n}{n-1}$, and
\[
 C_n=\binom{2n}{n}-\binom{2n}{n-1}.
\]
Double-count pairs consisting of an $n$-element subset of a $2n$-element
set and a distinguished element of that subset: there are
$n\binom{2n}{n}$ pairs. Removing the distinguished element instead gives
an $(n-1)$-element subset with one of its $n+1$ missing elements marked,
so $n\binom{2n}{n}=(n+1)\binom{2n}{n-1}$. Multiply the preceding
difference by $n+1$ to obtain the claimed identity. For $n=0$ it is
$C_0=1$. Subset enumeration, prefix reflection, and its inverse are
finite algorithms; their inverse identities and counting equations use
bounded induction. No real asymptotic estimate enters the proof.
\end{proof}

\begin{corollary}[Balanced parentheses are context-free and nonregular]
\label{count:dycknonreg}
The Dyck language is context-free but is not accepted by any finite DFA.
\end{corollary}
\begin{proof}
The grammar above proves context-freeness. Intersect $\mathcal D$ with
the regular language consisting of a block of opening parentheses followed
by a block of closing parentheses. The result is exactly
$\{\mathtt{(}^n\mathtt{)}^n:n\ge0\}$. A DFA for $\mathcal D$ would,
by the product construction, give a DFA for that language, contradicting
the regular pumping argument in Section~\ref{sec:regular}.
\end{proof}
